% Authored for ReadableSTEM's development set (dev-v2). Original text, no third-party material.
% Licence: MIT. Author: ReadableSTEM contributors.
\documentclass[11pt]{article}
% Computer Modern text fonts, as pdfLaTeX writes them; XeTeX defaults to Latin Modern.
\usepackage[OT1]{fontenc}
\renewcommand{\rmdefault}{cmr}
\renewcommand{\ttdefault}{cmtt}
\usepackage{amsmath,amssymb}
\usepackage[margin=1.15in]{geometry}
\usepackage{float,booktabs,listings}
\setcounter{secnumdepth}{-1}
\lstset{basicstyle=\ttfamily\small,columns=fixed,keepspaces=true,showstringspaces=false}
\setlength{\parindent}{0pt}\setlength{\parskip}{7pt}

\begin{document}

\begin{center}{\Large\bfseries Study Skills for Science: Notes, Examples and Key Terms}\end{center}

\section{Estimating Rates of Change from Tables of Measured Values Collected During Class Experiments}

A table of values is often the first form in which a rate of change is met. The difference quotient between neighbouring rows is the best guide to the behaviour of the underlying quantity.

\paragraph{Check the units before you start.} Convert every quantity to a single system of units before substituting any numbers.

\paragraph{Why does the sign change?} Dividing both sides of an inequality by a negative number reverses its direction.

\paragraph{Notation} Vectors are written in bold type and scalars in italics.

A short Python function for this step is shown below.

{\small\begin{verbatim}
def bisect(f, lo, hi, tol=1e-9):
    # f(lo) and f(hi) must have opposite signs
    while hi - lo > tol:
        mid = (lo + hi) / 2
        if f(lo) * f(mid) <= 0:
            hi = mid
        else:
            lo = mid
    return (lo + hi) / 2
\end{verbatim}}

A short Python function for this step is shown below.

\begin{lstlisting}[language=Python]
def trapezoid(f, a, b, n=100):
    h = (b - a) / n
    total = (f(a) + f(b)) / 2
    for k in range(1, n):
        total += f(a + k * h)
    return total * h
\end{lstlisting}

\section{Reading Graphs of Velocity Against Time and Recovering the Distance Travelled by an Object}

The area between a velocity graph and the time axis equals the displacement. Signs matter: regions below the axis count as motion in the opposite direction.

\textbf{Remember the order of operations:} Brackets first, then powers, then multiplication and division, and last addition and subtraction.

\paragraph{Do not round early!} Carry at least two more digits than you need until the last line of the calculation.

\paragraph{Caution} Calculators may be set to degrees or to radians; check before using trigonometric functions.

\noindent\begin{tabular}{@{}p{0.46\textwidth}@{\hspace{0.06\textwidth}}p{0.46\textwidth}@{}}
Reading list. Chapter 3 of the course notes, the practice sheet on rates of change, and the short article on measurement error. & Schedule. Lecture on Monday, laboratory on Wednesday, and a quiz at the start of class on Friday morning. \\
\end{tabular}

\subsubsection*{Example}

A cart rolls 12 m down a ramp in 4.0 s, starting from rest. Find its acceleration.

\subsubsection*{Solution}

Use d = (1/2) a t\textsuperscript{2}, so a = 2d / t\textsuperscript{2} = 2(12) / (4.0)\textsuperscript{2} = 1.5 m/s\textsuperscript{2}.

A short Python function for this step is shown below.

{\small\begin{verbatim}
def newton(f, df, x, steps=8):
    for _ in range(steps):
        x = x - f(x) / df(x)
    return x

print(newton(lambda t: t * t - 2, lambda t: 2 * t, 1.0))
\end{verbatim}}

A short Python function for this step is shown below.

\begin{lstlisting}[language=Python]
def gcd(a, b):
    while b:
        a, b = b, a % b
    return a

print(gcd(84, 36))
\end{lstlisting}

\section{How Rounding Errors Accumulate When a Long Calculation Is Carried Out on a Pocket Calculator}

Every rounded intermediate result introduces a small error, and these errors can grow. Keeping extra digits until the final step is a simple habit that avoids most of the damage.

\paragraph{Write down what you know.} A short list of given quantities and unknowns often shows which formula is needed.

\textbf{Is the answer reasonable?} Compare the size of the result with a rough estimate made in your head.

\paragraph{Reminder} Bring the data sheet and a ruler to the next lab session.

\begin{table}[H]
\centering
\caption{Measurements of a falling steel ball}
\begin{tabular}{lccc}
\toprule
\textbf{Trial} & \textbf{Height (m)} & \textbf{Time (s)} & \textbf{Speed (m/s)} \\
\midrule
1 & 0.50 & 0.32 & 3.1 \\
2 & 1.00 & 0.45 & 4.4 \\
3 & 1.50 & 0.55 & 5.4 \\
4 & 2.00 & 0.64 & 6.3 \\
5 & 2.50 & 0.71 & 7.0 \\
\bottomrule
\end{tabular}
\end{table}

A short Python function for this step is shown below.

{\small\begin{verbatim}
def sieve(n):
    flags = [True] * (n + 1)
    flags[0:2] = [False, False]
    for p in range(2, int(n ** 0.5) + 1):
        if flags[p]:
            for m in range(p * p, n + 1, p):
                flags[m] = False
    return [i for i, ok in enumerate(flags) if ok]
\end{verbatim}}

A short Python function for this step is shown below.

\begin{lstlisting}[language=Python]
def mean_and_variance(xs):
    n = len(xs)
    mean = sum(xs) / n
    var = sum((x - mean) ** 2 for x in xs) / (n - 1)
    return mean, var
\end{lstlisting}

\subsection*{Key Takeaways}
\begin{itemize}
\item A rate of change is a ratio of two differences.
\item Neighbouring table rows give the best estimate of a local rate.
\item State the unit of every rate you report.
\end{itemize}

\section{Choosing Between a Linear Model and an Exponential Model for a Small Set of Experimental Data}

A straight line and a curve that bends upward can both pass close to a handful of points. The choice should rest on the process that produced the data, not only on how well the curve fits.

\paragraph{Draw a picture first.} Even a crude sketch of the situation prevents many sign and direction errors.

\paragraph{State the assumptions clearly.} Ignoring air resistance or treating a gas as ideal is acceptable only if it is said aloud.

\paragraph{Convention} Angles are measured anticlockwise from the positive horizontal axis.

\subsubsection*{Example}

A solution is made by dissolving 18 g of salt in water to make 250 cm\textsuperscript{3}. Find the concentration in g per cm\textsuperscript{3}.

\subsubsection*{Solution}

The concentration is 18 / 250 = 0.072 g/cm\textsuperscript{3}.

A short Python function for this step is shown below.

{\small\begin{verbatim}
def fit_line(xs, ys):
    n = len(xs)
    sx, sy = sum(xs), sum(ys)
    sxy = sum(x * y for x, y in zip(xs, ys))
    sxx = sum(x * x for x in xs)
    slope = (n * sxy - sx * sy) / (n * sxx - sx ** 2)
    return slope, (sy - slope * sx) / n
\end{verbatim}}

A short Python function for this step is shown below.

\begin{lstlisting}[language=Python]
def transpose(m):
    return [list(row) for row in zip(*m)]

a = [[1, 2, 3], [4, 5, 6]]
print(transpose(a))
\end{lstlisting}

\subsection*{Check Your Understanding}
\begin{enumerate}
\item Explain in one sentence why the slope of a position graph is a velocity.
\item A graph is flat between t = 2 s and t = 5 s. What does that say about the motion?
\end{enumerate}

\section{A Careful Look at Units, Dimensions and the Common Mistakes Made When Converting Between Them}

A numerical answer without a unit is incomplete, and a unit with the wrong dimensions is a sign of an error earlier in the work. Checking both takes only a few seconds.

\textbf{Start from the definition.} When a formula is forgotten, the definition of the quantity usually rebuilds it.

\paragraph{Label every axis and every curve.} A graph without labels and units cannot be read by anyone else.

\noindent\begin{tabular}{@{}ll@{}}
Name: & \rule{7cm}{0.4pt} \\[10pt]
Student number: & \rule{7cm}{0.4pt} \\[10pt]
Date of laboratory: & \rule{7cm}{0.4pt} \\
\end{tabular}

A short Python function for this step is shown below.

{\small\begin{verbatim}
def binary_search(items, target):
    lo, hi = 0, len(items) - 1
    while lo <= hi:
        mid = (lo + hi) // 2
        if items[mid] == target:
            return mid
        if items[mid] < target:
            lo = mid + 1
        else:
            hi = mid - 1
    return -1
\end{verbatim}}

A short Python function for this step is shown below.

\begin{lstlisting}[language=Python]
counts = {}
for word in text.split():
    counts[word] = counts.get(word, 0) + 1
top = sorted(counts, key=counts.get, reverse=True)[:5]
\end{lstlisting}

\subsection*{Key Takeaways}
\begin{itemize}
\item Area under a velocity graph is displacement.
\item Regions below the axis count as negative.
\item Check the answer against a rough area estimate.
\end{itemize}

\section{From Words to Equations: Translating Short Problem Statements into Algebra Without Losing Information}

Most algebra mistakes are made before any algebra is done, when a sentence is turned into an equation. Writing down what every symbol stands for removes much of the risk.

\paragraph{One change at a time.} Vary a single quantity in each trial so that the effect can be traced to its cause.

\textbf{Keep a record of failed attempts?} They explain later why a particular approach was abandoned.

\begin{table}[H]
\centering
\caption{Selected properties of four metals}
\begin{tabular}{lccc}
\toprule
\textbf{Metal} & \textbf{Symbol} & \textbf{Density (g/cm\textsuperscript{3})} & \textbf{Melting point (K)} \\
\midrule
\textbf{Aluminium} & Al & 2.70 & 933 \\
\textbf{Copper} & Cu & 8.96 & 1358 \\
\textbf{Iron} & Fe & 7.87 & 1811 \\
\textbf{Zinc} & Zn & 7.14 & 693 \\
\bottomrule
\end{tabular}
\end{table}

\subsubsection*{Example}

The sum of two numbers is 25 and their difference is 7. Find both numbers.

\subsubsection*{Solution}

Adding the two equations gives twice the larger number, 32, so the larger number is 16 and the smaller is 9.

A short Python function for this step is shown below.

{\small\begin{verbatim}
class Vector:
    def __init__(self, x, y):
        self.x, self.y = x, y

    def dot(self, other):
        return self.x * other.x + self.y * other.y
\end{verbatim}}

A short Python function for this step is shown below.

\begin{lstlisting}[language=Python]
def celsius_to_kelvin(c):
    if c < -273.15:
        raise ValueError("below absolute zero")
    return c + 273.15
\end{lstlisting}

\subsection*{Check Your Understanding}
\begin{enumerate}
\item Why is it risky to round 1/3 to 0.33 in the first line of a long calculation?
\item How many significant figures does 0.00450 have?
\end{enumerate}

\section{Why Some Chemical Reactions Reach Equilibrium and What the Equilibrium Constant Tells Us About Them}

A reversible reaction does not stop; the forward and reverse rates become equal. The equilibrium constant compares the amounts of products and reactants once that balance has been reached.

\paragraph{Use consistent notation throughout the whole solution.} Switching between two symbols for the same quantity is a frequent source of confusion.

\paragraph{Look for a pattern in the first few cases.} A conjecture formed from small examples can then be tested on a larger one.

The same idea written in C looks like this.

{\small\begin{verbatim}
int sum_array(const int *a, int n) {
    int total = 0;
    for (int i = 0; i < n; i++) {
        total += a[i];
    }
    return total;
}
\end{verbatim}}

The same idea written in C looks like this.

\begin{lstlisting}[language=C]
void swap(int *p, int *q) {
    int tmp = *p;
    *p = *q;
    *q = tmp;
}
\end{lstlisting}

\subsection*{Key Takeaways}
\begin{itemize}
\item Round only at the end of a calculation.
\item Keep two spare digits in every intermediate step.
\item Report the answer with the precision of the data.
\end{itemize}

\section{Sketching Polynomial Functions by Finding Intercepts, Turning Points and the Long-Run Behaviour}

The shape of a polynomial can be predicted from a few facts: its degree, its roots, and the sign of the leading coefficient. A rough sketch based on these facts is often enough.

\textbf{Read the question twice.} The second reading usually reveals a condition that the first one skipped.

\paragraph{Substitute back to check.} Putting the solution into the original equation catches arithmetic slips and extraneous roots.

\noindent\begin{tabular}{@{}p{0.46\textwidth}@{\hspace{0.06\textwidth}}p{0.46\textwidth}@{}}
\rule{6cm}{0.4pt} & \rule{6cm}{0.4pt} \\[2pt]
Student signature & Instructor signature \\
\end{tabular}

\subsubsection*{Example}

A data set has the values 4.2, 5.1 and 6.0. Find the mean.

\subsubsection*{Solution}

The mean is (4.2 + 5.1 + 6.0) / 3 = 5.1.

The same idea written in C looks like this.

{\small\begin{verbatim}
struct point {
    double x;
    double y;
};

double norm2(struct point p) {
    return p.x * p.x + p.y * p.y;
}
\end{verbatim}}

The same idea written in C looks like this.

\begin{lstlisting}[language=C]
unsigned long factorial(unsigned n) {
    unsigned long r = 1;
    while (n > 1) {
        r *= n--;
    }
    return r;
}
\end{lstlisting}

\subsection*{Check Your Understanding}
\begin{enumerate}
\item Name one situation in which a straight line is the wrong model for data.
\item What does a negative residual tell you about a data point?
\end{enumerate}

\section{Interpreting the Slope and the Intercept of a Line of Best Fit in the Context of an Experiment}

The slope of a fitted line is a rate with units, and the intercept is a starting value. Reading both in the language of the experiment turns a formula into a statement about the world.

\paragraph{Estimate the uncertainty.} Half of the smallest scale division is a sensible starting point for a single reading.

\textbf{Take a break when stuck:} A short walk often clears the way to an idea that hours at the desk did not.

\begin{table}[H]
\centering
\caption{Summary of quiz scores by group}
\begin{tabular}{lccc}
\toprule
\textbf{Group} & \textbf{Students} & \textbf{Mean score} & \textbf{Standard deviation} \\
\midrule
\textbf{Morning} & 28 & 71.4 & 9.2 \\
\textbf{Afternoon} & 31 & 68.9 & 10.5 \\
\textbf{Evening} & 19 & 74.2 & 8.1 \\
\bottomrule
\end{tabular}
\end{table}

The same idea written in C looks like this.

{\small\begin{verbatim}
#include <stdio.h>

int main(void) {
    for (int k = 1; k <= 5; k++)
        printf("%d squared is %d\n", k, k * k);
    return 0;
}
\end{verbatim}}

The same idea written in C looks like this.

\begin{lstlisting}[language=C]
int count_bits(unsigned v) {
    int n = 0;
    for (; v; v >>= 1)
        n += v & 1u;
    return n;
}
\end{lstlisting}

\subsection*{Key Takeaways}
\begin{itemize}
\item Choose the model from the process, not from the plot alone.
\item Compare residuals, not only the correlation.
\item A good fit inside the data says little outside it.
\end{itemize}

\section{Sampling, Uncertainty and the Difference Between a Mistake in the Method and a Random Fluctuation}

No measurement is exact, and repeated readings scatter around a typical value. A systematic offset, by contrast, moves every reading in the same direction.

\paragraph{Explain the result in words.} A final sentence in plain language shows that the numbers have been understood.

\paragraph{Mind the difference between mass and weight.} Mass is measured in kilograms and weight, a force, in newtons.

\subsubsection*{Example}

A sample of 80 g of a radioactive isotope has a half-life of 6 hours. How much remains after 18 hours?

\subsubsection*{Solution}

Eighteen hours is 3 half-lives, so the mass is 80 / 2\textsuperscript{3} = 10 g.

The same idea written in C looks like this.

{\small\begin{verbatim}
double mean(const double *x, size_t n) {
    double s = 0.0;
    for (size_t i = 0; i < n; i++)
        s += x[i];
    return s / n;
}
\end{verbatim}}

From the command line the steps are as follows.

\begin{lstlisting}[language=bash]
for f in data/*.csv; do
    echo "$f has $(wc -l < "$f") lines"
done
\end{lstlisting}

\subsection*{Check Your Understanding}
\begin{enumerate}
\item A student writes 5 m + 3 s. What is wrong with this expression?
\item Convert 72 km/h to metres per second.
\end{enumerate}

\section{Counting Techniques for Arrangements and Selections, with Worked Examples from Games of Chance}

Counting problems become easier once it is clear whether order matters and whether repetition is allowed. The two questions fix which formula applies.

\textbf{Collect like terms before factoring.} Neater expressions have fewer chances for sign errors.

\paragraph{Compare with a limiting case.} Setting a parameter to zero or letting it grow often gives a result that is known already.

\noindent\begin{tabular}{@{}p{0.46\textwidth}@{\hspace{0.06\textwidth}}p{0.46\textwidth}@{}}
Given. The mass is 2.0 kg, the height is 5.0 m and the starting speed is zero. & Find. The speed of the mass just before it reaches the ground. \\
\end{tabular}

From the command line the steps are as follows.

{\small\begin{verbatim}
mkdir -p results/plots
python analyse.py --input raw.csv --out results/
ls -l results
\end{verbatim}}

From the command line the steps are as follows.

\begin{lstlisting}[language=bash]
cut -d, -f2 grades.csv | sort -n | uniq -c | sort -rn | head
\end{lstlisting}

\subsection*{Key Takeaways}
\begin{itemize}
\item Every equation must balance in its units.
\item Convert before you substitute.
\item Write the unit on every line of the work.
\end{itemize}

\section{Solving Systems of Linear Equations by Elimination and Recognising When No Solution Exists at All}

Elimination replaces a pair of equations by an equivalent pair in which one unknown has disappeared. A row that reduces to zero equals a nonzero number reveals that the system has no solution.

\paragraph{Report the final answer with its unit.} A number alone is not an answer in any science course.

\textbf{Share the work with a partner?} Explaining a step aloud is the quickest way to find a gap in it.

\begin{table}[H]
\centering
\caption{Half-lives of some radioactive isotopes}
\begin{tabular}{lccc}
\toprule
\textbf{Isotope} & \textbf{Protons} & \textbf{Neutrons} & \textbf{Half-life} \\
\midrule
\textbf{\textsuperscript{14}C} & 6 & 8 & 5730 years \\
\textbf{\textsuperscript{60}Co} & 27 & 33 & 5.27 years \\
\textbf{\textsuperscript{131}I} & 53 & 78 & 8.0 days \\
\textbf{\textsuperscript{238}U} & 92 & 146 & 4.5 billion years \\
\bottomrule
\end{tabular}
\end{table}

\subsubsection*{Example}

A rectangle has a perimeter of 30 cm and one side of 4 cm. Find its area.

\subsubsection*{Solution}

The other side is 30/2 - 4 = 11 cm, so the area is 44 cm\textsuperscript{2}.

From the command line the steps are as follows.

{\small\begin{verbatim}
find . -name "*.log" -mtime +30 -print
grep -c "ERROR" run.log
\end{verbatim}}

From the command line the steps are as follows.

\begin{lstlisting}[language=bash]
tar -czf backup.tar.gz notes/ figures/
tar -tzf backup.tar.gz | head -n 5
\end{lstlisting}

\subsection*{Check Your Understanding}
\begin{enumerate}
\item State two assumptions behind the ideal gas law.
\item What quantity stays constant at equilibrium: the amounts or the rates?
\end{enumerate}

\section{Writing a Clear Laboratory Report That Another Student Could Follow Without Asking Any Questions}

A good report states the question, the method, the data and the conclusion, in that order. Each part should be short enough to read in a minute.

From the command line the steps are as follows.

{\small\begin{verbatim}
name="lab3"
echo "Preparing $name"
cp template.tex "$name.tex"
head -n 3 "$name.tex" 
\end{verbatim}}

A database query for the same summary is given next.

\begin{lstlisting}[language=SQL]
SELECT student, AVG(score) AS mean_score
FROM quiz
GROUP BY student
HAVING COUNT(*) >= 3
ORDER BY mean_score DESC;
\end{lstlisting}

\section{Checking an Answer for Reasonableness Before Submitting It: Estimates, Limiting Cases and Symmetry}

Before an answer is handed in, a quick estimate, a limiting case or a symmetry argument can reveal a slip. Such checks are cheap and surprisingly effective.

A database query for the same summary is given next.

{\small\begin{verbatim}
SELECT s.name, c.title
FROM enrolment AS e
JOIN student AS s ON s.id = e.student_id
JOIN course AS c ON c.id = e.course_id;
\end{verbatim}}

In JavaScript the idea is expressed like this.

{\small\begin{verbatim}
function clamp(x, lo, hi) {
  return Math.min(hi, Math.max(lo, x));
}

console.log(clamp(12, 0, 10));
\end{verbatim}}

\section{Key Terms}

\begin{description}

\item[Derivative] The limit of the ratio of the change in a function to the change in its input; it gives the slope of the tangent line.

\item[Integral] A limit of sums that measures accumulated change, such as the area under a curve.

\item[Limit] The value that a function approaches as its input gets arbitrarily close to a chosen point.

\item[Matrix] A rectangular array of numbers arranged in rows and columns.

\item[Determinant] A single number computed from a square matrix that tells whether the matrix can be inverted.

\item[Eigenvalue] A scalar by which a matrix stretches a special direction called an eigenvector.

\item[Vector] A quantity with both a size and a direction, often written as a column of numbers.

\item[Scalar] A quantity described by a single number and, in physics, a unit.

\item[Domain] The set of all inputs for which a function is defined.

\item[Range] The set of all outputs that a function actually produces.

\item[Asymptote] A line that a curve approaches without ever reaching it.

\item[Polynomial] A sum of terms, each a constant times a whole-number power of a variable.

\item[Logarithm] The exponent to which a base must be raised to give a stated number.

\item[Sequence] An ordered list of numbers, usually given by a rule.

\item[Series] The sum of the terms of a sequence.

\item[Probability] A number between zero and one that measures how likely an event is.

\item[Mean] The sum of a set of values divided by how many values there are.

\item[Variance] The average squared distance of the values in a data set from their mean.

\item[Hypothesis] A statement about the world that an experiment is designed to test.

\item[Isotope] One of several forms of an element that have the same number of protons but different numbers of neutrons.

\item[Ion] An atom or molecule that carries a net electric charge.

\item[Mole] The amount of substance that contains as many particles as there are atoms in twelve grams of carbon-12.

\item[Molarity] The number of moles of solute in one litre of solution.

\item[Catalyst] A substance that speeds up a reaction without being used up.

\item[Equilibrium] A state in which the forward and reverse rates of a process are equal and the amounts stop changing.

\item[Half-life] The time needed for half of the nuclei in a radioactive sample to decay.

\item[Velocity] The rate of change of position, with a direction.

\item[Acceleration] The rate of change of velocity.

\item[Momentum] The product of an object's mass and its velocity.

\item[Frequency] The number of complete cycles of a wave that pass a point each second.

\item[Wavelength] The distance between two matching points on neighbouring cycles of a wave.

\item[Amplitude] The largest displacement of a wave from its rest position.

\item[Algorithm] A finite list of precise steps that solves a class of problems.

\item[Variable] A name that stands for a value which may change.

\end{description}

\end{document}
